% !TEX root = ../高等概率论.tex
% =====================================================================
%  分布：分布与分布函数、分位数变换、离散/有密度/奇异分布、耦合、二进制展开、正态尾部估计
% =====================================================================
\lecture{随机变量的分布与分布函数，分布函数的刻画与分位数变换，离散分布、有密度的分布与 Cantor 分布，耦合与随机序，二进制展开与均匀分布，正态分布的尾部估计}
\label{lec:distributions}

随机变量 $X$ 本身是 $\Omega$ 上的函数，而 $\Omega$ 往往是抽象的、不必关心的。我们真正关心的是 $X$ 取各种值的概率，这由 $X$ 在 $\R$ 上诱导的测度，即 $X$ 的分布来描述。本讲说明：$\R$ 上的概率分布与分布函数一一对应，并且每个分布函数都可以由 $(0,1)$ 上的均匀分布经过一次变换得到。后半部分讨论几个具体的例子，它们展示了分布可以有多么不同的形态，以及“把几个随机变量放在同一个概率空间上”（耦合）这一想法的威力。

\topic{分布与分布函数}

\begin{summary}
随机变量把 $P$ “推”到 $\R$ 上，得到分布。分布由分布函数决定，分布函数由单调性、右连续性与两端的极限刻画。
\end{summary}

\begin{definition}[分布]\label{def:distributions:law}
设 $(\Omega,\cF,P)$ 是概率空间，$X:(\Omega,\cF)\to(S,\cS)$ 可测。称 $(S,\cS)$ 上的函数
\[
\mu_X(B)=P(X\in B)=P(X^{-1}(B)),\qquad B\in\cS
\]
为 $X$ 的\term{分布}（distribution）或\term{律}（law），也记作 $P\circ X^{-1}$，称为 $P$ 在 $X$ 下的\term{像测度}（image measure, pushforward）。若两个随机元 $X,Y$（可以定义在不同的概率空间上）有相同的分布，称它们\term{同分布}（equal in distribution），记作 $X\eqd Y$。
\end{definition}

$\mu_X$ 是 $(S,\cS)$ 上的概率测度：$\mu_X(S)=P(\Omega)=1$；若 $B_n\in\cS$ 两两不交，则事件 $\{X\in B_n\}$ 两两不交（$X(\omega)$ 不能同时落在两个不交的集合中），且 $\{X\in\bigcup_nB_n\}=\bigcup_n\{X\in B_n\}$，于是 $\mu_X(\bigcup_nB_n)=\sum_nP(X\in B_n)=\sum_n\mu_X(B_n)$。

\begin{example}[分布不决定随机变量]
在 $([0,1],\cR\cap[0,1],\lambda)$ 上令 $X(\omega)=\omega$，$Y(\omega)=1-\omega$。对 $0\leqslant x\leqslant1$，$P(X\leqslant x)=\lambda([0,x])=x$，$P(Y\leqslant x)=\lambda([1-x,1])=x$。于是 $X\eqd Y$（由推论~\ref{cor:measures:distribution-function}），但除了 $\omega=1/2$ 之外处处 $X(\omega)\neq Y(\omega)$。同分布只是关于“随机规律”的陈述，与两个函数是否相等无关。
\end{example}

\begin{definition}[分布函数]\label{def:distributions:df}
随机变量 $X$ 的\term{分布函数}（distribution function）是 $F_X(x)=P(X\leqslant x)=\mu_X((-\infty,x])$，$x\in\R$。
\end{definition}

记 $F(x-)=\lim_{y\uparrow x}F(y)$ 为左极限。

\begin{theorem}[分布函数的性质]\label{thm:distributions:df-properties}
设 $F$ 是随机变量 $X$ 的分布函数。则
\begin{enumerate}
\item $F$ 单调不减；
\item $\lim_{x\to\infty}F(x)=1$，$\lim_{x\to-\infty}F(x)=0$；
\item $F$ 右连续；
\item $F(x-)=P(X<x)$；
\item $P(X=x)=F(x)-F(x-)$。
\end{enumerate}
于是 $P(a<X\leqslant b)=F(b)-F(a)$，$P(a\leqslant X\leqslant b)=F(b)-F(a-)$，等等。
\end{theorem}
\begin{proof}
(1) 若 $x\leqslant y$，则 $\{X\leqslant x\}\subseteq\{X\leqslant y\}$，由单调性（定理~\ref{thm:measures:basic}(2)）即得。

(2) 由 (1)，两个极限存在，只需沿 $x=\pm n$ 计算。$\{X\leqslant n\}\uparrow\Omega$（每个 $X(\omega)$ 是实数，终将不超过 $n$），由从下连续，$F(n)\to P(\Omega)=1$。$\{X\leqslant-n\}\downarrow\varnothing$（没有实数不超过一切 $-n$），由从上连续（概率测度有限），$F(-n)\to0$。

(3) 由 (1)，只需对 $y_n=x+1/n$ 验证。$\{X\leqslant x+1/n\}\downarrow\{X\leqslant x\}$，由从上连续，$F(x+1/n)\to F(x)$。

(4) 同理只需对 $y_n=x-1/n$ 验证：$\{X\leqslant x-1/n\}\uparrow\{X<x\}$，由从下连续即得。

(5) $\{X=x\}=\{X\leqslant x\}\setminus\{X<x\}$，由定理~\ref{thm:measures:basic}(3) 与 (4)。
\end{proof}

\begin{proposition}\label{prop:distributions:jumps}
单调不减函数 $F:\R\to\R$ 至多有可数个间断点。
\end{proposition}
\begin{proof}
单调函数在每点有左、右极限，且 $F(x-)\leqslant F(x+)$；$F$ 在 $x$ 不连续当且仅当 $F(x-)<F(x+)$。对每个间断点 $x$，取有理数 $q(x)\in(F(x-),F(x+))$。若 $x<y$ 都是间断点，取 $z\in(x,y)$，则 $F(x+)\leqslant F(z)\leqslant F(y-)$，所以区间 $(F(x-),F(x+))$ 与 $(F(y-),F(y+))$ 不交，$q(x)\neq q(y)$。于是 $x\mapsto q(x)$ 是从间断点集到 $\Q$ 的单射。
\end{proof}

由推论~\ref{cor:measures:distribution-function}，$X\eqd Y$ 当且仅当 $F_X=F_Y$。下面的定理说明，定理~\ref{thm:distributions:df-properties} 的 (1)(2)(3) 已经完全刻画了分布函数。

\begin{theorem}[分布函数的刻画]\label{thm:distributions:characterization}
设 $F:\R\to\R$ 满足定理~\ref{thm:distributions:df-properties} 的 (1)(2)(3)。则存在概率空间上的随机变量 $X$，使 $F_X=F$。
\end{theorem}

我们给出两个证明。

\begin{proof}[证明一（用 Lebesgue--Stieltjes 测度）]
$F$ 是 Stieltjes 测度函数，由定理~\ref{thm:extension:lebesgue-stieltjes}，存在 $(\R,\cR)$ 上的测度 $\mu_F$，使 $\mu_F((a,b])=F(b)-F(a)$。由从下连续，$\mu_F(\R)=\lim_n\mu_F((-n,n])=\lim_n(F(n)-F(-n))=1$，且 $\mu_F((-\infty,x])=\lim_n\mu_F((-n,x])=F(x)-0=F(x)$。取概率空间 $(\R,\cR,\mu_F)$，$X(\omega)=\omega$，则 $P(X\leqslant x)=\mu_F((-\infty,x])=F(x)$。
\end{proof}

第二个证明更有用，它把一切分布都归结为均匀分布。

\begin{definition}[分位数函数]\label{def:distributions:quantile}
设 $F$ 满足定理~\ref{thm:distributions:df-properties} 的 (1)(2)(3)。对 $u\in(0,1)$，令
\[
F^{-1}(u)=\inf\{x\in\R: F(x)\geqslant u\},
\]
称为 $F$ 的\term{分位数函数}（quantile function）或\term{广义逆}（generalized inverse）。
\end{definition}

\begin{lemma}\label{lem:distributions:quantile}
设 $u\in(0,1)$。则 $F^{-1}(u)$ 是有限实数，并且对一切 $x\in\R$，
\[
F^{-1}(u)\leqslant x\quad\Longleftrightarrow\quad u\leqslant F(x).
\]
特别地 $F^{-1}$ 单调不减。
\end{lemma}
\begin{proof}
令 $I_u=\{x:F(x)\geqslant u\}$。由于 $F(x)\to1>u$（$x\to\infty$），$I_u$ 非空；由于 $F(x)\to0<u$（$x\to-\infty$），存在 $M$ 使 $x<M$ 时 $F(x)<u$，于是 $I_u$ 有下界 $M$。所以 $F^{-1}(u)=\inf I_u$ 是有限实数。

由单调性，若 $x\in I_u$ 且 $y>x$，则 $y\in I_u$；所以 $I_u$ 是一个向右无界的区间，端点为 $a=F^{-1}(u)$。$a$ 本身属于 $I_u$：取 $x_n\in I_u$，$x_n\downarrow a$（若 $a\in I_u$ 已成立则无需此步），由右连续性，$F(a)=\lim_nF(x_n)\geqslant u$。所以 $I_u=[a,\infty)$，即 $F(x)\geqslant u\iff x\geqslant F^{-1}(u)$，这就是要证的等价式。

若 $u<v$，则 $I_v\subseteq I_u$，下确界 $F^{-1}(u)\leqslant F^{-1}(v)$。
\end{proof}

\begin{proof}[定理~\ref{thm:distributions:characterization} 的证明二（分位数变换）]
取概率空间 $((0,1),\cR\cap(0,1),\lambda)$，令 $X(\omega)=F^{-1}(\omega)$。由引理~\ref{lem:distributions:quantile}，$X$ 单调不减，从而是 Borel 函数（定理~\ref{thm:random-variables:composition}(3) 的证明同样适用于 $(0,1)$ 上的单调函数），并且
\[
\{\omega\in(0,1): X(\omega)\leqslant x\}=\{\omega\in(0,1):\omega\leqslant F(x)\}=(0,F(x)]\cap(0,1).
\]
由于 $0\leqslant F(x)\leqslant1$，这个集合的 Lebesgue 测度是 $F(x)$（$F(x)=1$ 时它是 $(0,1)$，测度仍为 $1$）。所以 $P(X\leqslant x)=F(x)$。
\end{proof}

\begin{remark}[分位数变换]
证明二说明：若 $U$ 是\emph{任意}概率空间上服从 $(0,1)$ 上均匀分布的随机变量，即 $P(U\leqslant u)=u$（$0\leqslant u\leqslant1$），且 $U\in(0,1)$，则 $F^{-1}(U)$ 的分布函数是 $F$。计算机生成随机数正是这样做的：先生成均匀随机数，再作分位数变换。当 $F$ 连续且严格递增时，$F^{-1}$ 就是通常的反函数；一般情形下，$F$ 的平台（常值段）对应 $F^{-1}$ 的跳跃，$F$ 的跳跃对应 $F^{-1}$ 的平台。
\end{remark}

\begin{corollary}
$\R$ 上的概率测度、满足 (1)(2)(3) 的函数 $F$、以及随机变量的同分布类，三者一一对应。
\end{corollary}
\begin{proof}
概率测度 $\mu\mapsto F(x)=\mu((-\infty,x])$ 满足 (1)(2)(3)（与定理~\ref{thm:distributions:df-properties} 的证明相同，取 $X$ 为恒等映射）；由推论~\ref{cor:measures:distribution-function} 这是单射；由证明一，它是满射。随机变量 $X\mapsto\mu_X$ 把同分布类映为概率测度，由定理~\ref{thm:distributions:characterization}，每个概率测度都是某个随机变量的分布。
\end{proof}

\topic{离散分布、有密度的分布与 Cantor 分布}

\begin{summary}
分布函数可以有跳跃，可以是可微函数的积分，也可以连续却“几乎处处导数为零”。三种情形分别对应离散分布、有密度的分布与奇异连续分布。
\end{summary}

\begin{definition}[离散分布]\label{def:distributions:discrete}
若存在可数集 $D\subseteq\R$ 使 $\mu(D)=1$，称 $\R$ 上的概率测度 $\mu$（或其分布函数）为\term{离散的}（discrete）。
\end{definition}

设 $D=\{x_1,x_2,\dots\}$，$p_i=\mu(\{x_i\})$，则 $\mu=\sum_ip_i\delta_{x_i}$，$F(x)=\sum_ip_i\one_{[x_i,\infty)}(x)$。初等概率论中的 Bernoulli 分布、二项分布、几何分布、Poisson 分布都是离散分布。离散分布函数在 $x_i$ 处跳跃 $p_i$，但它不一定是“阶梯函数”：

\begin{example}[跳跃点稠密的离散分布]\label{ex:distributions:dense-jumps}
设 $q_1,q_2,\dots$ 是 $\Q$ 的一个排列，$\alpha_i>0$，$\sum_i\alpha_i=1$（例如 $\alpha_i=2^{-i}$）。令
\[
F(x)=\sum_{i=1}^\infty\alpha_i\one_{[q_i,\infty)}(x).
\]
由\egref{ex:extension:LS-examples} 的 (4)，$F$ 是分布函数，对应的分布 $\sum_i\alpha_i\delta_{q_i}$ 集中在 $\Q$ 上，是离散的。它具有以下看似矛盾的性质：
\begin{enumerate}
\item $F$ 严格递增：若 $x<y$，取有理数 $q_j\in(x,y)$，则 $F(y)-F(x)=\sum_{i:x<q_i\leqslant y}\alpha_i\geqslant\alpha_j>0$。
\item $F$ 在每个有理数 $q_j$ 处有跳跃 $\alpha_j$，从而间断点稠密：由定理~\ref{thm:distributions:df-properties}(5)，$F(q_j)-F(q_j-)=\mu(\{q_j\})=\alpha_j$。
\item $F$ 在每个无理数处连续：$\mu(\{x\})=0$。
\end{enumerate}
图像上，$F$ 在任何一个小区间里都有无穷多个跳跃，却没有任何一段是平的。它提醒我们：“离散分布”指的是测度集中在可数集上，而不是分布函数长得像楼梯。
\end{example}

\begin{definition}[密度]\label{def:distributions:density}
若存在非负函数 $f$ 使
\[
F(x)=\int_{-\infty}^xf(y)\dd y\qquad\text{对一切}\ x\in\R,
\]
称分布函数 $F$（或其分布）有\term{密度}（density）$f$。
\end{definition}

在建立 Lebesgue 积分之前，这里的积分可以理解为（广义）Riemann 积分，$f$ 取为分段连续函数即可。有密度的分布函数是连续的，所以 $P(X=x)=0$ 对一切 $x$ 成立。反之，给定非负、积分为 $1$ 的 $f$，上式定义的 $F$ 满足 (1)(2)(3)，从而是分布函数。

\begin{example}[三个重要的密度]\label{ex:distributions:densities}
\begin{enumerate}
\item \term{$(0,1)$ 上的均匀分布}：$f(x)=\one_{(0,1)}(x)$；$F(x)=0$（$x\leqslant0$），$F(x)=x$（$0\leqslant x\leqslant1$），$F(x)=1$（$x>1$）。
\item \term{参数为 $\lambda>0$ 的指数分布}（exponential distribution）：$f(x)=\lambda e^{-\lambda x}\one_{[0,\infty)}(x)$；$F(x)=(1-e^{-\lambda x})\one_{[0,\infty)}(x)$。它的分位数函数是 $F^{-1}(u)=-\log(1-u)/\lambda$，所以若 $U$ 是均匀分布的，$-\log(1-U)/\lambda$ 服从指数分布。
\item \term{标准正态分布}（standard normal distribution）：$\varphi(x)=(2\pi)^{-1/2}e^{-x^2/2}$。由数学分析中熟知的 $\int_{-\infty}^\infty e^{-x^2/2}\dd x=\sqrt{2\pi}$，$\varphi$ 是密度。它的分布函数记作 $\Phi$，没有初等表达式；本讲最后给出它的尾部估计。
\end{enumerate}
\end{example}

最后是一个既无跳跃又无密度的例子。

\begin{example}[Cantor 分布]\label{ex:distributions:cantor}
回忆 Cantor 集 $C=\bigcap_nC_n$（习题~\exref{ex:extension:A4}）：$C_n$ 是 $2^n$ 个长度为 $3^{-n}$ 的闭区间的不交并，$C_{n+1}$ 由 $C_n$ 的每个区间去掉中间的开三分之一得到。$\lambda(C_n)=(2/3)^n$，所以 $\lambda(C)=0$。

\emph{构造。}令 $F_n:\R\to[0,1]$ 为这样的连续函数：$x\leqslant0$ 时 $F_n(x)=0$，$x\geqslant1$ 时 $F_n(x)=1$；在 $C_n$ 的每个区间上线性地增加 $2^{-n}$；在 $[0,1]\setminus C_n$ 的每个区间上为常数。（从左到右，第 $j$ 个区间 $J$ 上从 $(j-1)2^{-n}$ 增加到 $j2^{-n}$。）$F_n$ 单调不减。

\emph{$F_n$ 一致收敛。}设 $J=[s,s+3^{-n}]$ 是 $C_n$ 的第 $j$ 个区间。$C_{n+1}$ 把 $J$ 分成左、右两个长度为 $3^{-n-1}$ 的区间，中间去掉 $(s+3^{-n-1},s+2\cdot3^{-n-1})$。$F_{n+1}$ 在 $J$ 的两个端点处取值与 $F_n$ 相同（分别是 $(j-1)2^{-n}$ 与 $j2^{-n}$，因为 $J$ 左侧 $C_{n+1}$ 的区间个数是 $2(j-1)$），在 $J$ 的中间段上取值 $(j-1)2^{-n}+2^{-n-1}$。所以在 $J$ 上，$F_n$ 与 $F_{n+1}$ 都取值于 $[(j-1)2^{-n},j2^{-n}]$，差不超过 $2^{-n}$；在 $C_n$ 之外二者相等。于是 $\sup_x\abs{F_{n+1}(x)-F_n(x)}\leqslant2^{-n}$，$\{F_n\}$ 一致收敛到一个函数 $F$。

\emph{$F$ 的性质。}一致收敛的连续函数列的极限连续，单调不减函数列的极限单调不减，所以 $F$ 连续、单调不减，$F=0$ 在 $(-\infty,0]$，$F=1$ 在 $[1,\infty)$，从而 $F$ 是分布函数。$F$ 称为 \term{Cantor 函数}，对应的分布 $\mu$ 称为 \term{Cantor 分布}。若 $(a,b)$ 是某个 $[0,1]\setminus C_n$ 中的开区间（去掉的区间），则对一切 $m\geqslant n$，$F_m$ 在 $[a,b]$ 上是常数，故 $F$ 在 $[a,b]$ 上是常数。

\emph{$\mu$ 集中在 $C$ 上。}设 $(a,b)$ 是一个去掉的开区间。由 $F$ 在 $[a,b]$ 上为常数，$\mu((a,b))=\lim_k\mu((a,b-1/k])=\lim_k(F(b-1/k)-F(a))=0$。$[0,1]\setminus C_n$ 是有限个去掉的开区间的并，所以 $\mu([0,1]\setminus C_n)=0$；又 $\mu((-\infty,0))=F(0-)=0$，$\mu((1,\infty))=1-F(1)=0$。于是 $\mu(C_n)=1$，$C_n\downarrow C$，由从上连续 $\mu(C)=1$。

\emph{结论。}$\mu$ 没有原子（$F$ 连续，$\mu(\{x\})=0$），所以不是离散的，也不含离散部分；另一方面 $\mu(C)=1$ 而 $\lambda(C)=0$，$\mu$ 与 Lebesgue 测度“生活在不交的集合上”。这样的分布称为\term{奇异的}（singular）。它不可能有密度：学了 Lebesgue 积分以后，若 $f$ 是密度，则 $\mu(C)=\int_Cf\dd\lambda=0$（在零测集上的积分为零），矛盾。直观上，$F$ 在 $C$ 之外（一个测度为 $1$ 的开集上）导数为零，却从 $0$ 连续地爬升到了 $1$，全部的增长都发生在一个长度为零的集合上。
\end{example}

\begin{remark}
以后会证明（Lebesgue 分解定理）：$\R$ 上的每个概率测度都可以唯一地写成 $\mu=a\mu_d+b\mu_{ac}+c\mu_s$，其中 $a+b+c=1$，$\mu_d$ 离散，$\mu_{ac}$ 有密度，$\mu_s$ 连续且奇异。本讲作业 D 证明其中较容易的一半：分出离散部分。
\end{remark}

\topic{耦合与随机序}

\begin{summary}
比较两个分布时，一个非常有效的方法是：把两个随机变量构造在同一个概率空间上，使它们之间满足某种逐点关系（比如 $X\leqslant Y$），然后读出分布之间的关系。这就是耦合。
\end{summary}

\begin{definition}[耦合]\label{def:distributions:coupling}
设 $\mu,\nu$ 是 $(\R,\cR)$ 上的概率测度。若 $(X,Y)$ 是某个概率空间上的随机向量，$X$ 的分布是 $\mu$，$Y$ 的分布是 $\nu$，则称 $(X,Y)$ 是 $\mu$ 与 $\nu$ 的一个\term{耦合}（coupling）。
\end{definition}

耦合总是存在的。例如设 $F,G$ 是 $\mu,\nu$ 的分布函数，由定理~\ref{thm:extension:Rd}，$(\R^2,\cR^2)$ 上存在测度 $\pi$ 使 $\pi((a_1,b_1]\times(a_2,b_2])=(F(b_1)-F(a_1))(G(b_2)-G(a_2))$；取 $X,Y$ 为两个坐标投影，则 $P(X\leqslant x)=\pi((-\infty,x]\times\R)=F(x)$，$P(Y\leqslant y)=G(y)$（用从下连续取极限）。这是“独立”的耦合。耦合的艺术在于选择一个\emph{有用的}耦合。

\begin{example}[两枚硬币]\label{ex:distributions:two-coins}
设 $0\leqslant p\leqslant q\leqslant1$。硬币甲正面朝上的概率是 $p$，硬币乙是 $q$。各独立地掷 $n$ 次，记正面次数分别为 $S$ 与 $T$。直觉告诉我们，对每个 $k$，$P(S\geqslant k)\leqslant P(T\geqslant k)$。直接比较
\[
\sum_{j\geqslant k}\binom njp^j(1-p)^{n-j}\quad\text{与}\quad\sum_{j\geqslant k}\binom njq^j(1-q)^{n-j}
\]
并不容易。耦合的证明如下。

取概率空间 $\Omega=(0,1]^n$，$\cF=\cR^n\cap(0,1]^n$，$P=\lambda_n$（$n$ 维 Lebesgue 测度限制在 $(0,1]^n$ 上，由定理~\ref{thm:extension:Rd} 它在矩形上取值为各边长之积，从而 $P(\Omega)=1$）。记 $U_i(\omega)=\omega_i$，令
\[
X_i=\one_{\{U_i\leqslant p\}},\qquad Y_i=\one_{\{U_i\leqslant q\}},\qquad S=\sum_{i=1}^nX_i,\qquad T=\sum_{i=1}^nY_i.
\]
\emph{分布是对的。}对 $a=(a_1,\dots,a_n)\in\{0,1\}^n$，$\{X_1=a_1,\dots,X_n=a_n\}=I_1\times\cdots\times I_n$，其中 $a_i=1$ 时 $I_i=(0,p]$，$a_i=0$ 时 $I_i=(p,1]$。这是一个矩形，所以
\[
P(X_1=a_1,\dots,X_n=a_n)=\prod_{i=1}^np^{a_i}(1-p)^{1-a_i}=p^{\sum a_i}(1-p)^{n-\sum a_i}.
\]
这正是“$n$ 次独立的、正面概率为 $p$ 的抛掷”的概率模型。对 $j$ 求和，$\{S=j\}$ 是 $\binom nj$ 个这样的事件的不交并，$P(S=j)=\binom njp^j(1-p)^{n-j}$，即 $S$ 服从二项分布 $\mathrm{Bin}(n,p)$。同理 $T\sim\mathrm{Bin}(n,q)$。

\emph{逐点比较。}由 $p\leqslant q$，$\{U_i\leqslant p\}\subseteq\{U_i\leqslant q\}$，故 $X_i\leqslant Y_i$ 处处成立，$S\leqslant T$ 处处成立。于是 $\{S\geqslant k\}\subseteq\{T\geqslant k\}$，
\[
P(S\geqslant k)\leqslant P(T\geqslant k).
\]
由于这两个概率只依赖于 $S$、$T$ 的分布，我们证明了关于二项分布的不等式 $P(\mathrm{Bin}(n,p)\geqslant k)\leqslant P(\mathrm{Bin}(n,q)\geqslant k)$，而证明中一个求和都没有算。

这里的想法是：用同一个 $U_i$ 同时“驱动”两枚硬币。$U_i\leqslant p$ 时两枚都是正面，$p<U_i\leqslant q$ 时只有乙是正面，$U_i>q$ 时两枚都是反面。
\end{example}

两枚硬币的例子是下面这个一般结论的特例。

\begin{definition}[随机序]\label{def:distributions:stochastic-order}
设 $X,Y$ 是随机变量（可以在不同的概率空间上），分布函数为 $F,G$。若 $F(x)\geqslant G(x)$ 对一切 $x$ 成立，即 $P(X>x)\leqslant P(Y>x)$ 对一切 $x$ 成立，称 $X$ \term{随机地小于} $Y$（stochastically smaller），记作 $X\preceq_{\mathrm{st}}Y$。
\end{definition}

\begin{theorem}[随机序的耦合刻画]\label{thm:distributions:stochastic-coupling}
设 $F,G$ 是分布函数。则 $F(x)\geqslant G(x)$ 对一切 $x$ 成立，当且仅当存在同一概率空间上的随机变量 $X,Y$，$F_X=F$，$F_Y=G$，且 $X\leqslant Y$ 处处成立。
\end{theorem}
\begin{proof}
充分性：若 $X\leqslant Y$，则 $\{Y\leqslant x\}\subseteq\{X\leqslant x\}$，所以 $G(x)=P(Y\leqslant x)\leqslant P(X\leqslant x)=F(x)$。

必要性：设 $F\geqslant G$。在 $((0,1),\cR\cap(0,1),\lambda)$ 上令 $X(\omega)=F^{-1}(\omega)$，$Y(\omega)=G^{-1}(\omega)$。由定理~\ref{thm:distributions:characterization} 证明二，$F_X=F$，$F_Y=G$。对 $u\in(0,1)$，若 $G(x)\geqslant u$，则 $F(x)\geqslant G(x)\geqslant u$，所以 $\{x:G(x)\geqslant u\}\subseteq\{x:F(x)\geqslant u\}$，取下确界得 $F^{-1}(u)\leqslant G^{-1}(u)$，即 $X\leqslant Y$ 处处成立。
\end{proof}

在两枚硬币的例子中，$n=1$ 时 $X_1$ 与 $Y_1$ 恰好是这个分位数耦合（差一个 $\omega\mapsto1-\omega$ 的翻转）。

\begin{historynote}
耦合方法今天是概率论的基本工具。它的一个惊人的新近应用是：2026 年初，A.\ Song 在一维情形下证明，存在普适常数 $c>0$，使得对任何满足 $\abs X\leqslant1$、均值为 $0$ 的随机变量 $X$，都可以在同一个概率空间上构造三个服从标准正态分布的随机变量 $G_1,G_2,G_3$（它们不必独立），使 $cX=G_1+G_2+G_3$。随后 D.\ M.\ Hua、A.\ Song 与 S.\ Tudose 把结论推广到 $\R^n$ 中的随机向量，并由此解决了 M.\ Talagrand 关于 Gauss 空间中“大集合的 Minkowski 和包含凸体”的凸性问题。论文还指出，一般情形下只用两个标准正态随机向量是不够的。这里的要点正是耦合：三个 $G_i$ 各自的分布都是固定的，自由度全在于如何把它们“摆放”在同一个概率空间上。
\end{historynote}

\topic{二进制展开与均匀分布}

\begin{summary}
均匀分布与“无穷次公平抛硬币”是同一个对象的两种看法：均匀随机数的二进制各位是公平硬币；反过来，把公平硬币的结果当作二进制各位，就得到均匀随机数。
\end{summary}

对 $\omega\in[0,1)$ 与 $n\geqslant1$，令
\[
\xi_n(\omega)=\lfloor2^n\omega\rfloor-2\lfloor2^{n-1}\omega\rfloor .
\]

\begin{proposition}[均匀随机数的二进制各位]\label{prop:distributions:binary-digits}
在 $([0,1),\cR\cap[0,1),\lambda)$ 上：
\begin{enumerate}
\item 每个 $\xi_n$ 取值于 $\{0,1\}$，是随机变量，并且 $\omega=\sum_{n=1}^\infty\xi_n(\omega)2^{-n}$；
\item 对任意 $n\geqslant1$ 与 $a_1,\dots,a_n\in\{0,1\}$，
\[
\{\xi_1=a_1,\dots,\xi_n=a_n\}=\Bigl[\frac k{2^n},\frac{k+1}{2^n}\Bigr),\qquad k=\sum_{i=1}^na_i2^{n-i},
\]
从而 $P(\xi_1=a_1,\dots,\xi_n=a_n)=2^{-n}$。
\end{enumerate}
\end{proposition}
\begin{proof}
记 $t=2^{n-1}\omega$，$m=\lfloor t\rfloor$，则 $m\leqslant t<m+1$，$2m\leqslant2t<2m+2$，所以 $\lfloor2t\rfloor\in\{2m,2m+1\}$，即 $\xi_n\in\{0,1\}$。$\omega\mapsto\lfloor2^n\omega\rfloor$ 单调，是 Borel 函数，所以 $\xi_n$ 是随机变量。由 $\xi_i$ 的定义逐项相消，
\[
\sum_{i=1}^n\xi_i2^{n-i}=\sum_{i=1}^n\bigl(\lfloor2^i\omega\rfloor2^{n-i}-\lfloor2^{i-1}\omega\rfloor2^{n-i+1}\bigr)=\lfloor2^n\omega\rfloor-\lfloor\omega\rfloor=\lfloor2^n\omega\rfloor,
\]
这里用了 $\omega\in[0,1)$ 时 $\lfloor\omega\rfloor=0$。于是 $\sum_{i\leqslant n}\xi_i2^{-i}=2^{-n}\lfloor2^n\omega\rfloor\in(\omega-2^{-n},\omega]$，令 $n\to\infty$ 得 (1)。

(2) 由上式，$\xi_1=a_1,\dots,\xi_n=a_n$ 蕴含 $\lfloor2^n\omega\rfloor=k$。反之，若 $\lfloor2^n\omega\rfloor=k$，则对 $j\leqslant n$，$\lfloor2^j\omega\rfloor=\lfloor2^{j-n}\lfloor2^n\omega\rfloor\rfloor=\lfloor k/2^{n-j}\rfloor$（这用到对正整数 $N$ 与实数 $s$，$\lfloor s/N\rfloor=\lfloor\lfloor s\rfloor/N\rfloor$），这由 $k$ 决定，从而 $\xi_1,\dots,\xi_n$ 由 $k$ 决定；而 $k=\sum_ia_i2^{n-i}$ 的二进制表示是唯一的，所以 $\xi_i=a_i$。于是事件等于 $\{\lfloor2^n\omega\rfloor=k\}=[k2^{-n},(k+1)2^{-n})$，其 Lebesgue 测度为 $2^{-n}$（命题~\ref{prop:extension:lebesgue}(1)）。
\end{proof}

这个命题有两层意义。第一，它\emph{构造}了“无穷次公平抛硬币”的概率模型：$\xi_1,\xi_2,\dots$ 定义在同一个概率空间上，任何有限次的结果都是等可能的。第二，下面的定理说明，这个性质反过来刻画了均匀分布，与概率空间的具体选择无关。

\begin{theorem}[公平硬币给出均匀分布]\label{thm:distributions:coins-to-uniform}
设 $(\Omega,\cF,P)$ 是概率空间，$\xi_1,\xi_2,\dots:\Omega\to\{0,1\}$ 是随机变量，并且对任意 $n$ 与 $a_1,\dots,a_n\in\{0,1\}$，
\begin{equation}\label{eq:distributions:fair}
P(\xi_1=a_1,\dots,\xi_n=a_n)=2^{-n}.
\end{equation}
令 $U=\sum_{n=1}^\infty\xi_n2^{-n}$。则 $U$ 是随机变量，并且 $U$ 服从 $[0,1]$ 上的均匀分布：$P(U\leqslant x)=x$ 对一切 $x\in[0,1]$ 成立。
\end{theorem}
\begin{proof}
\emph{可测性。}部分和 $U_n=\sum_{i\leqslant n}\xi_i2^{-i}$ 是随机变量，$U_n\uparrow U$ 处处成立（级数非负，且 $U\leqslant\sum_n2^{-n}=1$），由定理~\ref{thm:random-variables:limits}，$U$ 是取值于 $[0,1]$ 的随机变量。

\emph{二进有理点处的估计。}固定 $n$。对 $0\leqslant j<2^n$，令 $a^{(j)}\in\{0,1\}^n$ 为 $j$ 的 $n$ 位二进制表示，$E_j=\{(\xi_1,\dots,\xi_n)=a^{(j)}\}$。$E_0,\dots,E_{2^n-1}$ 两两不交，并为 $\Omega$，由 \eqref{eq:distributions:fair}，$P(E_j)=2^{-n}$。在 $E_j$ 上，$U_n=j2^{-n}$，而 $0\leqslant U-U_n\leqslant\sum_{i>n}2^{-i}=2^{-n}$，所以
\[
E_j\subseteq\{j2^{-n}\leqslant U\leqslant(j+1)2^{-n}\}.
\]
现在设 $d=k2^{-n}$，$0\leqslant k\leqslant2^n$。若 $j\geqslant k$，则在 $E_j$ 上 $U\geqslant j2^{-n}\geqslant d$；所以 $\{U<d\}\subseteq\bigcup_{j<k}E_j$。若 $j<k$，则在 $E_j$ 上 $U\leqslant(j+1)2^{-n}\leqslant d$；所以 $\bigcup_{j<k}E_j\subseteq\{U\leqslant d\}$。取概率，
\begin{equation}\label{eq:distributions:dyadic}
P(U<d)\leqslant k2^{-n}=d\leqslant P(U\leqslant d).
\end{equation}

\emph{一般的 $x$。}设 $0<x<1$。由于二进有理数在 $[0,1]$ 中稠密，对充分大的 $m$ 可以取二进有理数 $d_m,d_m'\in[0,1]$，使 $x-2^{-m}<d_m'<x<d_m<x+2^{-m}$。由 \eqref{eq:distributions:dyadic} 与单调性，
\[
d_m'\leqslant P(U\leqslant d_m')\leqslant P(U\leqslant x)\leqslant P(U<d_m)\leqslant d_m.
\]
令 $m\to\infty$ 得 $P(U\leqslant x)=x$。$x=1$ 时 $U\leqslant1$ 处处成立；$x=0$ 时，$P(U\leqslant0)\leqslant P(U<2^{-m})\leqslant2^{-m}$ 对一切 $m$ 成立，故为 $0$。
\end{proof}

\begin{remark}
条件 \eqref{eq:distributions:fair} 就是以后要定义的“$\xi_1,\xi_2,\dots$ 独立同分布，$P(\xi_n=0)=P(\xi_n=1)=1/2$”。定理的证明只用到 \eqref{eq:distributions:fair} 这个关于有限维分布的条件，这与上一讲 $C[0,1]$ 的例子是同一个道理：有限维分布决定了一切。作为推论，$P(\text{从某项起}\ \xi_n\ \text{全为}\ 1)=0$：该事件包含于 $\bigcup_N\bigcap_{n\geqslant N}\{\xi_n=1\}$，而 $P(\xi_N=\cdots=\xi_{N+m}=1)=2^{-m-1}$。所以“二进制展开的歧义”只发生在一个零概率事件上。
\end{remark}

\topic{正态分布的尾部估计}

\begin{summary}
标准正态分布函数没有初等表达式，但它的尾部概率 $P(X\geqslant x)$ 可以被两个初等函数紧紧夹住：它大约是 $\varphi(x)/x$。
\end{summary}

设 $X$ 服从标准正态分布，即对 Borel 集 $B$，$P(X\in B)=\frac1{\sqrt{2\pi}}\int_Be^{-x^2/2}\dd x$。记 $\varphi(x)=(2\pi)^{-1/2}e^{-x^2/2}$，$\overline\Phi(x)=P(X\geqslant x)=\int_x^\infty\varphi(y)\dd y$。

\begin{theorem}[Gauss 尾部估计]\label{thm:distributions:gaussian-tail}
对一切 $x>0$，
\[
\Bigl(\frac1x-\frac1{x^3}\Bigr)\varphi(x)\ \leqslant\ \overline\Phi(x)\ \leqslant\ \frac1x\varphi(x),
\qquad\text{并且}\qquad \overline\Phi(x)\leqslant\frac12e^{-x^2/2}.
\]
特别地，$x\to\infty$ 时 $\overline\Phi(x)\big/\bigl(\varphi(x)/x\bigr)\to1$。
\end{theorem}
\begin{proof}
\emph{上界。}当 $y\geqslant x>0$ 时 $y/x\geqslant1$，而 $\varphi'(y)=-y\varphi(y)$，所以
\[
\overline\Phi(x)=\int_x^\infty\varphi(y)\dd y\leqslant\int_x^\infty\frac yx\varphi(y)\dd y=\frac1x\bigl[-\varphi(y)\bigr]_{y=x}^{y=\infty}=\frac{\varphi(x)}x .
\]

\emph{下界。}令 $h(y)=(y^{-1}-y^{-3})\varphi(y)$，$y>0$。求导，
\[
h'(y)=(-y^{-2}+3y^{-4})\varphi(y)+(y^{-1}-y^{-3})(-y)\varphi(y)=-(1-3y^{-4})\varphi(y).
\]
由于 $h(y)\to0$（$y\to\infty$），
\[
(x^{-1}-x^{-3})\varphi(x)=h(x)=\int_x^\infty(1-3y^{-4})\varphi(y)\dd y\leqslant\int_x^\infty\varphi(y)\dd y=\overline\Phi(x).
\]

\emph{第二个上界。}作变换 $y=x+z$，并利用 $e^{-(x+z)^2/2}=e^{-x^2/2}e^{-xz}e^{-z^2/2}\leqslant e^{-x^2/2}e^{-z^2/2}$（$x,z\geqslant0$），
\[
\overline\Phi(x)=\int_0^\infty\varphi(x+z)\dd z\leqslant e^{-x^2/2}\int_0^\infty\varphi(z)\dd z=\frac12e^{-x^2/2},
\]
最后一步用了 $\varphi$ 关于原点对称，积分为 $1$。

\emph{渐近。}由两个界，$1-x^{-2}\leqslant\overline\Phi(x)\big/(\varphi(x)/x)\leqslant1$，令 $x\to\infty$。
\end{proof}

\begin{example}
取 $x=3$。$\varphi(3)=e^{-4.5}/\sqrt{2\pi}\approx0.004432$，定理给出 $0.001313\leqslant P(X\geqslant3)\leqslant0.001477$，而真实值约为 $0.001350$。第二个上界给出 $\frac12e^{-4.5}\approx0.005554$，粗糙得多，但它在 $x$ 较小时也有效，并且形式简单：例如 $x=0$ 时它给出精确值 $1/2$，而 $\varphi(x)/x$ 在 $x\to0$ 时趋于无穷。两类界各有用处：前者刻画精确的渐近，后者（以及它的变形）在大偏差与集中不等式中经常出现。
\end{example}

% =====================================================================
\homework{Cantor 分布的概率构造，最大耦合与全变差距离，分布的离散部分与连续部分}[Cantor 分布，最大耦合，离散部分与连续部分]
\label{hw:distributions}

\exercisepart{基本习题}

\begin{exercises}{A}{课堂内容的补充}
\item\label{ex:distributions:A1} 设 $F$ 是分布函数，$F^{-1}$ 是分位数函数。证明：
\begin{questions}
\item $F^{-1}$ 在 $(0,1)$ 上左连续；
\item $F(F^{-1}(u))\geqslant u$，等号在 $F$ 于 $F^{-1}(u)$ 处连续时成立；
\item $F^{-1}(F(x))\leqslant x$（当 $0<F(x)<1$ 时），并举例说明不等号可以严格；
\item $F$ 连续时 $F(F^{-1}(u))=u$；$F$ 严格递增时 $F^{-1}(F(x))=x$。
\end{questions}
\item 设 $X$ 的分布函数 $F$ 连续。证明 $Y=F(X)$ 服从 $(0,1)$ 上的均匀分布：$P(Y\leqslant y)=y$，$y\in[0,1]$。若 $F$ 不连续，结论是否成立？
\item 设 $X$ 有连续密度 $f$，$P(\alpha<X<\beta)=1$，$g$ 在 $(\alpha,\beta)$ 上严格递增且可微。证明 $g(X)$ 在 $(g(\alpha),g(\beta))$ 上有密度 $f(g^{-1}(y))/g'(g^{-1}(y))$。$g(x)=ax+b$（$a>0$）时，密度是什么？
\item 设 $X$ 服从标准正态分布。求 $e^X$ 的密度（\term{对数正态分布}）与 $X^2$ 的密度（自由度为 $1$ 的 \term{$\chi^2$ 分布}）。
\item 用定理~\ref{thm:distributions:gaussian-tail} 给出 $P(X\geqslant4)$ 的上下界，其中 $X$ 服从标准正态分布。
\item 设 $X$ 服从参数为 $\lambda$ 的指数分布。证明\term{无记忆性}：$P(X>s+t\mid X>s)=P(X>t)$，$s,t\geqslant0$。又证明 $\lceil X\rceil$ 服从\term{几何分布}：$P(\lceil X\rceil=k)=(1-p)^{k-1}p$，$k\geqslant1$，并求 $p$。
\item 设 $X\eqd Y$，$f:\R\to\R$ 是 Borel 函数。证明 $f(X)\eqd f(Y)$。
\item 设 $X$ 是随机变量，$F$ 是它的分布函数。求 $\max\{X,0\}$、$\abs X$、$-X$ 的分布函数（用 $F$ 表示）。
\item 验证\egref{ex:distributions:dense-jumps} 中 $F$ 的右连续性：对 $\varepsilon>0$，取 $N$ 使 $\sum_{i>N}\alpha_i<\varepsilon$，再利用有限和 $\sum_{i\leqslant N}\alpha_i\one_{[q_i,\infty)}$ 的右连续性。
\item 在\egref{ex:distributions:two-coins} 中，不用耦合，直接证明 $n=2$，$k=1$ 与 $k=2$ 时的不等式。再证明：对 $k\geqslant1$，$p\mapsto P(\mathrm{Bin}(n,p)\geqslant k)$ 的导数为 $n\binom{n-1}{k-1}p^{k-1}(1-p)^{n-k}\geqslant0$，从而得到另一个证明。
\end{exercises}

\begin{probleme}{B}{Cantor 分布的概率构造}
设 $(\Omega,\cF,P)$ 上的 $\xi_1,\xi_2,\dots$ 满足 \eqref{eq:distributions:fair}。令 $X=\sum_{n=1}^\infty2\xi_n3^{-n}$。本题证明 $X$ 服从 Cantor 分布。记 $F$ 为 Cantor 函数，$C_n$、$C$ 如\egref{ex:distributions:cantor}。
\begin{questions}
\item 证明 $X$ 是随机变量，取值于 $C$。（提示：前 $n$ 位决定了 $X$ 落在 $C_n$ 的哪一个区间中。）
\item 设 $J$ 是 $C_n$ 从左数第 $j$ 个区间。证明 $\{X\in J\}=\{(\xi_1,\dots,\xi_n)=b^{(j-1)}\}$，其中 $b^{(j-1)}$ 是 $j-1$ 的 $n$ 位二进制表示。从而 $P(X\in J)=2^{-n}$。
\item 设 $J=[s,s+3^{-n}]$ 是 $C_n$ 的第 $j$ 个区间。证明 $P(X\leqslant s)=(j-1)2^{-n}$，$P(X\leqslant s+3^{-n})=j2^{-n}$，并与 $F$ 在这两点的值比较。
\item 由 B3) 与 $F_X$、$F$ 的单调性，证明 $\abs{F_X(x)-F(x)}\leqslant2^{-n}$ 对一切 $x$ 成立，从而 $F_X=F$。
\item 令 $X'=\sum_{n\geqslant1}2\xi_{n+1}3^{-n}$，则 $X=\frac23\xi_1+\frac13X'$，且 $\xi_2,\xi_3,\dots$ 也满足 \eqref{eq:distributions:fair}。仿照 B2)--B4)，证明对 $x\in[0,1]$，$P(\xi_1=0,\ X'\leqslant x)=\frac12F(x)$，$P(\xi_1=1,\ X'\leqslant x)=\frac12F(x)$。由此推出 Cantor 函数的\term{自相似性}：$F(x/3)=\frac12F(x)$，$F(\frac23+\frac x3)=\frac12+\frac12F(x)$。
\item 设 $x=\sum_n\varepsilon_n3^{-n}\in C$，$\varepsilon_n\in\{0,2\}$。证明 $\{X\leqslant x\}$ 与“$(\xi_1,\xi_2,\dots)$ 按字典序不超过 $(\varepsilon_1/2,\varepsilon_2/2,\dots)$”这一事件相等，再由定理~\ref{thm:distributions:coins-to-uniform} 证明 $F(x)=\sum_n(\varepsilon_n/2)2^{-n}$。这与第~\ref{lec:random-variables} 讲思考题 E 中 Cantor 函数的定义一致。
\end{questions}
\end{probleme}

\begin{probleme}{C}{最大耦合与全变差距离}
设 $\mu,\nu$ 是可数集 $D\subseteq\R$ 上的概率分布，记 $\mu(x)=\mu(\{x\})$，$\nu(x)=\nu(\{x\})$。定义\term{全变差距离}（total variation distance）
\[
d_{TV}(\mu,\nu)=\sup_{A\subseteq D}\abs{\mu(A)-\nu(A)}.
\]
本题证明 \term{耦合不等式}（coupling inequality）及其等号成立条件：$d_{TV}(\mu,\nu)=\min\{P(X\neq Y):(X,Y)\ \text{是}\ \mu,\nu\ \text{的耦合}\}$。
\begin{questions}
\item 证明 $d_{TV}(\mu,\nu)=\frac12\sum_{x\in D}\abs{\mu(x)-\nu(x)}=1-\sum_{x\in D}\min\{\mu(x),\nu(x)\}$，并且上确界在 $A=\{x:\mu(x)>\nu(x)\}$ 处取到。
\item 设 $(X,Y)$ 是任意耦合。证明对任意 $A$，$\mu(A)-\nu(A)=P(X\in A)-P(Y\in A)\leqslant P(X\in A,Y\notin A)\leqslant P(X\neq Y)$。从而 $d_{TV}(\mu,\nu)\leqslant P(X\neq Y)$。
\item 记 $m(x)=\min\{\mu(x),\nu(x)\}$，$\delta=d_{TV}(\mu,\nu)$，设 $0<\delta<1$。在可数集 $D\times D$ 上定义
\[
\pi(x,y)=m(x)\one_{\{x=y\}}+\frac{(\mu(x)-m(x))(\nu(y)-m(y))}{\delta}.
\]
证明 $\pi$ 是 $D\times D$ 上的概率分布，它的两个边缘分布分别是 $\mu$ 与 $\nu$。（提示：$\sum_y(\nu(y)-m(y))=\delta$。）
\item 取离散概率空间 $(D\times D,2^{D\times D},\pi)$ 与坐标投影 $(X,Y)$。证明 $P(X\neq Y)\leqslant\delta$，从而等号成立。（提示：$\mu(x)-m(x)$ 与 $\nu(x)-m(x)$ 不能同时为正。）处理 $\delta=0$ 与 $\delta=1$ 的情形。
\item 应用：设 $\mu=\mathrm{Bernoulli}(p)$，$\nu=\mathrm{Poisson}(p)$（$\nu(k)=e^{-p}p^k/k!$）。证明 $d_{TV}(\mu,\nu)=p(1-e^{-p})\leqslant p^2$。这是以后 Poisson 逼近的出发点。
\end{questions}
\end{probleme}

\begin{probleme}{D}{分布函数的离散部分与连续部分}
设 $F$ 是分布函数，$D=\{x:F(x)>F(x-)\}$ 是它的间断点集（由命题~\ref{prop:distributions:jumps} 可数），跳跃为 $p(x)=F(x)-F(x-)$。
\begin{questions}
\item 证明对任意 $a<b$，$\sum_{x\in D\cap(a,b]}p(x)\leqslant F(b)-F(a)$。（提示：先对有限个间断点证明。）从而 $\sum_{x\in D}p(x)\leqslant1$。
\item 令 $F_d(x)=\sum_{y\in D,\,y\leqslant x}p(y)$，$F_c=F-F_d$。证明 $F_d$ 单调不减、右连续，并且在每个 $x\in D$ 处跳跃 $p(x)$，在 $D$ 之外连续。
\item 由 D1) 证明 $F_c$ 单调不减。
\item 证明 $F_c$ 连续。（提示：计算 $F_c(x)-F_c(x-)$。）
\item 设 $a=\sum_{x\in D}p(x)$。证明：若 $0<a<1$，则 $F=aG_d+(1-a)G_c$，其中 $G_d=F_d/a$ 是离散分布函数，$G_c=F_c/(1-a)$ 是连续分布函数，并且这样的分解是唯一的。
\item 用 D5) 与\egref{ex:distributions:cantor} 构造一个分布函数，它既有跳跃，也有奇异连续的部分，还有有密度的部分。
\end{questions}
\end{probleme}

\thinkingpart

\begin{probleme}[问题]{E}{Mills 比的更精确估计}
设 $x>0$，记 $R(x)=\overline\Phi(x)/\varphi(x)$，称为 \term{Mills 比}（Mills ratio）。
\begin{questions}
\item 证明 $R'(x)=xR(x)-1$。
\item 由 E1) 与定理~\ref{thm:distributions:gaussian-tail} 证明 $R(x)\geqslant x/(1+x^2)$。（提示：考虑 $g(x)=R(x)-x/(1+x^2)$，证明 $g(\infty)=0$ 且 $g'<0$。）并与 $1/x-1/x^3$ 比较：哪个下界在 $x$ 较小时更好？
\item 反复分部积分，证明对每个 $k\geqslant0$，
\[
\sum_{j=0}^{2k+1}\frac{(-1)^j(2j-1)!!}{x^{2j+1}}\leqslant R(x)\leqslant\sum_{j=0}^{2k}\frac{(-1)^j(2j-1)!!}{x^{2j+1}},
\]
其中 $(-1)!!=1$。这个级数对固定的 $x$ 发散，却给出越来越精确的渐近展开。
\end{questions}
\end{probleme}

\begin{probleme}*[问题]{F}{分布函数收敛与分位数收敛}
设 $F_n,F$ 是分布函数，并且在 $F$ 的每个连续点 $x$ 处 $F_n(x)\to F(x)$。
\begin{questions}
\item 设 $u\in(0,1)$，$x<F^{-1}(u)$ 是 $F$ 的连续点。证明 $F(x)<u$，从而 $n$ 充分大时 $F_n(x)<u$，$F_n^{-1}(u)\geqslant x$。由此 $\liminf_nF_n^{-1}(u)\geqslant F^{-1}(u)$。
\item 设 $v>u$，$y>F^{-1}(v)$ 是 $F$ 的连续点。证明 $\limsup_nF_n^{-1}(u)\leqslant y$，从而 $\limsup_nF_n^{-1}(u)\leqslant F^{-1}(u+)$。
\item 证明：在 $F^{-1}$ 的每个连续点 $u$ 处，$F_n^{-1}(u)\to F^{-1}(u)$。
\item 由 F3) 与“单调函数至多可数个间断点”，在 $((0,1),\lambda)$ 上构造随机变量 $X_n\eqd F_n$，$X\eqd F$，使 $X_n\to X$ a.s.。（这是以后 Skorokhod 表示定理的一维情形。）
\end{questions}
\end{probleme}

\exercisepart{课外知识：一个均匀随机数包含无穷多个均匀随机数}

\begin{probleme}{G}{分拆二进制各位}
在 $([0,1),\lambda)$ 上，$\xi_n$ 是二进制各位（命题~\ref{prop:distributions:binary-digits}）。取一个双射 $\N\times\N\to\N$，$(k,m)\mapsto n(k,m)$，令 $U_k=\sum_{m\geqslant1}\xi_{n(k,m)}2^{-m}$。
\begin{questions}
\item 设 $N$ 是正整数，$b_1,\dots,b_N\in\{0,1\}$，$i_1<\cdots<i_N$ 是正整数。证明 $P(\xi_{i_1}=b_1,\dots,\xi_{i_N}=b_N)=2^{-N}$。（提示：把事件写成 $\{\xi_1,\dots,\xi_{i_N}\}$ 的若干取值组合的不交并，数一数组合的个数。）
\item 证明对每个 $k$，$\{\xi_{n(k,m)}\}_{m\geqslant1}$ 满足 \eqref{eq:distributions:fair}，从而 $U_k$ 服从均匀分布。
\item 设 $x_1,\dots,x_K$ 都是形如 $j2^{-M}$ 的二进有理数。证明 $P(U_1<x_1,\dots,U_K<x_K)=x_1\cdots x_K$。
\item 由 G3) 与逼近论证证明：对一切 $x_1,\dots,x_K\in[0,1]$，$P(U_1\leqslant x_1,\dots,U_K\leqslant x_K)=x_1\cdots x_K$。
\item 结合分位数变换，说明：给定任意一列分布函数 $G_1,G_2,\dots$，可以在 $([0,1),\lambda)$ 上构造随机变量 $Y_k=G_k^{-1}(U_k)$，使对任意 $K$ 与 $y_1,\dots,y_K$，$P(Y_1\leqslant y_1,\dots,Y_K\leqslant y_K)=\prod_kG_k(y_k)$。以后我们会说：$Y_1,Y_2,\dots$ 是独立的，分布分别为 $G_k$。（注意 $U_k$ 可能等于 $0$，这时 $G_k^{-1}(0)$ 需要另外约定，但这只发生在零概率事件上。）
\end{questions}
\end{probleme}
